\SourceLeaf{43}care nullus dubitaverit. Legem ibi, licet non summa generalitate expressam,
signis concinnioribus circumscribere, et ex genuino fonte rigida demonstratione
derivare, operae pretium mihi visum est.

Praemonenda autem ante omnia quaedam, de novo signo in exhibendis serierum
terminis generalibus utilissimo, brevibus exponam. In seriebus nimirum,
imprimis complicatioribus saepe accidit, ut aliam termini impares, aliam
pares legem sequi videantur, et tunc in binas, a se invicem seiunctas dirimi
soleant. Quod, cum non sine prolixitate fieri possit et, ut ita dicam,
unitatem expressionis, qua sola tractabilis redditur series indefinita,
tollat, nisi aliqua ratione medela adferatur, serierum usum arctis
circumscribet limitibus. Sponte autem ex methodo, qua simplicissima quae
dari potest inter terminos seriei pares et impares discrepantia, ut nimirum
alteri positivi, alteri negativi sint, in ipsum terminum generalem infertur,
generalior ista, quam nunc proponere conor, sese offerre videtur. Sit
terminus seriei alicuius generalis $m^{\text{tus}} = Q$; si illum positivum
velis existente $m$ pari, negativum $m$ impari, hoc praefixo factore
$(-1)^m$; sin contrario ordine procedere signa malis, praeposito
$(-1)^{m+1}$ exhibeo\HistoricalNote{*)}{*)}{Adsensum huic signo non recusare dignatus est, cum ipsi de illo
proponerem, Ill.\ \textsc{Kaestnerus}. Qua auctoritate fretus saepius iam in
disquisitionibus analyticis idem illud adhibui.}. Simili ratione, sit expressio generalis termini
$m^{\text{ti}}$ existente $m$ pari $= P$, existente impari $= Q$, ita, ut
$P$ aliam prorsus, quam $Q$ legem sequatur. Dico fore generatim terminum
$m^{\text{tum}}$ $\frac{1-(-1)^m}{2}Q+\frac{1-(-1)^{m+1}}{2}P$. Ex hac enim
expressione, si fuerit $m$ numerus par, ob $(-1)^m=1$, evadente
$\frac{1-(-1)^m}{2}=0$, et ob $(-1)^{m+1}=-1$, evadente
$\frac{1-(-1)^{m+1}}{2}=1$, eliditur pars prima, fitque altera uti fas est
$=P$. Sin vero $m$ fuerit impar, cum sit $(-1)^{m+1}=1$, adeoque
$\frac{1-(-1)^{m+1}}{2}=0$, et simul $(-1)^m=-1$ et inde
$\frac{1-(-1)^m}{2}=1$, evanescente parte posteriore restat prior ut
assumta fuit $=Q$. Ut unicum tantum exemplum adferam, ex quo notae utilitas
elucescat, sumam, quod infra demonstrabitur, posse potestatem indefinitam
cosinus dati cuiusdam anguli v.\,g.\ $\cos\varphi^n$ semper in seriem
secundum cosinus multiplorum eiusdem anguli expandi, ea tamen lege, ut ista,
existente $n$ pari, secundum multipla paria procedat, sitque adeo terminus
illius $r^{\text{tus}}$, misso coefficiente, $\cos 2r\varphi$; existente
autem $n$ impari secundum multipla imparia, sitque tunc terminus
eiusdem $r^{\text{tus}} = \cos(2r+1)\varphi$. Tunc expressio generalis sic
adornari potest: fore seriei in quam $\cos\varphi^n$ evol%
